% TEMPLATE-MILC-v3.tex - plantilla maestra UNICA de las guias MILC v3.
%
% La usan las cuatro familias de guias, cada una con su driver:
%   · Grados 6-11          -> scripts/build-guias-g11.py
%   · Bebras               -> scripts/build-guias-bebras.py
%   · Semillero            -> scripts/build-guias-semillero.py
%   · Territorio interior  -> scripts/build-guias-territorio-interior.py
%
% Antes habia tres copias casi identicas (~400 lineas iguales, 6 distintas):
% cada mejora habia que aplicarla tres veces, y en la practica no se hacia
% (el motor de recursos vivia solo en dos de ellas). Lo que varia entre
% familias esta parametrizado en 6 placeholders que cada driver llena
% (se nombran aqui SIN su sintaxis real: el builder reemplaza por texto plano,
% tambien dentro de los comentarios, y un valor multilinea rompe el preambulo):
%
%   HEADER_IZQ        encabezado izquierdo de pagina
%   HEADER_DER        encabezado derecho de pagina
%   PORTADA_ETIQUETA  rotulo sobre el numero grande (GRADO/MOMENTO/MODULO)
%   PORTADA_SERIE     linea de serie de la portada
%   PORTADA_ID        identificador de la guia en la portada
%   PORTADA_CIRCULO   el circulito de grado (°), o vacio si no aplica
%   PDF_TITULO        titulo en texto plano (sin \textbf ni comillas TeX) para
%                     los metadatos del PDF (/Title); lo produce lib_xelatex.texto_pdf
%
% Composicion (auditoria 2026-09-03): las bandas de estacion piden espacio
% (\Needspace*) y prohiben el salto justo despues (\nopagebreak), asi no quedan
% huerfanas al pie; las cajas largas (softbox, drawbox, infoband, puentebox) son
% `breakable`, asi una caja de media pagina ya no arrastra todo a la siguiente
% dejando la anterior al 40 %. Todo texto sobre mostaza o verde va en milcNegro
% (blanco sobre #E5B400 da 1,93:1 y sobre #58A923 2,95:1; ninguno pasa WCAG AA).
% El resto de claves las documenta content/guias/_SCHEMA.md. El script valida
% que no queden placeholders sin reemplazar antes de escribir el archivo final.

% Etiquetado PDF (latex-lab): /Lang, estructura y texto alternativo de las
% figuras, para lectores de pantalla. Debe ir ANTES de \documentclass.
\DocumentMetadata{lang=es-CO,tagging=on}
\documentclass[11pt,letterpaper]{article}
% headheight=24pt: la cabecera derecha lleva dos lineas (identidad de la guia
% y estacion actual). headsep se acorta en la misma medida para que el bloque
% de texto quede exactamente donde estaba con headheight=12pt.
\usepackage[margin=2.0cm,headheight=24pt,headsep=13pt]{geometry}
\usepackage[table]{xcolor}
\usepackage{fontspec}
\usepackage[spanish]{babel}
\usepackage{tikz}
% Librerías TikZ para los diagramas de las guías (bloque `recursos:`):
%  · babel  → babel en español vuelve activos caracteres como «>», y TikZ los usa
%             en sus opciones (>=stealth, ->). Sin esta librería, cualquier
%             diagrama con flechas revienta con «\language@active@arg> has an
%             extra }». Debe ir ANTES de que se dibuje nada.
%  · positioning        → sintaxis «below left=2pt and 3pt».
%  · decorations.pathreplacing → llaves ({brace}) para señalar rangos.
\usetikzlibrary{babel,positioning,decorations.pathreplacing}
\usepackage{graphicx}
% Circuitos: el trabajo de electrónica se hace hoy con micro:bit y sus sensores
% integrados (MakeCode, sin cableado), así que no se necesitan esquemáticos.
% Si algún día entran montajes con protoboard, basta descomentar esta línea:
% los diagramas .tex/.tikz declarados en `recursos:` ya se incrustan solos.
% \usepackage{circuitikz}
\usepackage[most]{tcolorbox}
\usepackage{tabularx}
\usepackage{array}
\usepackage{fancyhdr}
\usepackage{titlesec}
\usepackage[document]{ragged2e}  % bandera a la derecha en todo el cuerpo
\usepackage{enumitem}
\usepackage{microtype}
\usepackage{etoolbox}
\usepackage{needspace}
% hyperref al final del preambulo: metadatos del PDF (titulo, autor, idioma,
% familia), marcadores por estacion (\pdfbookmark) y enlaces sin recuadro.
\usepackage[hidelinks,
  pdftitle={Concebir tu libro — género, tema, audiencia y escaleta},
  pdfauthor={PhD. Álvaro Cárdenas Orozco},
  pdfsubject={Tecnología e Informática · Grado 10 · Guía 2 · MILC v3},
  pdflang={es-CO},
  bookmarksopen=true,bookmarksnumbered=false]{hyperref}

% Helvetica Neue no trae la flecha U+2192, asi que cada «→» escrito literalmente
% en el YAML se compilaba a NADA, en silencio (462 flechas en 61 guias: rutas del
% tipo «paleta → tipografia → lockup» salian rotas). Mapeamos el caracter a su
% version matematica, que si tiene glifo. Asi el autor puede seguir escribiendo
% «→» comodamente en el YAML.
\usepackage{newunicodechar}
\newunicodechar{→}{\ensuremath{\rightarrow}}
% Mismo problema con los simbolos de marca y vinetas: el «✓» de «una palomita ✓»
% salia como cuadrito vacio en el PDF de todas las guias que lo usaban.
\usepackage{pifont}
\newunicodechar{✓}{\ding{51}}
\newunicodechar{✗}{\ding{55}}
\newunicodechar{✔}{\ding{52}}
\newunicodechar{•}{\textbullet}
\newunicodechar{≥}{\ensuremath{\geq}}
\newunicodechar{≤}{\ensuremath{\leq}}

\setmainfont{Helvetica Neue}
\setsansfont{Helvetica Neue}
\setlength{\parindent}{0pt}
\setlength{\parskip}{6pt}
% Sin particion de palabras: en 6.º-7.º una palabra cortada por guion al final
% de linea («Comuni-cacion») frena la lectura mucho mas que una linea corta.
\hyphenpenalty=10000
\exhyphenpenalty=10000
\pretolerance=2000
\tolerance=3000
\setlength{\emergencystretch}{3em}
\linespread{1.10}
\renewcommand{\arraystretch}{1.25}
\setlist{leftmargin=5mm,itemsep=1mm,topsep=1mm}
\newcolumntype{Y}{>{\RaggedRight\arraybackslash}X}

\definecolor{milcMagenta}{HTML}{E6007E}
\definecolor{milcVino}{HTML}{5A0038}
\definecolor{milcTurquesa}{HTML}{0093A5}
\definecolor{milcVerde}{HTML}{58A923}
\definecolor{milcMostaza}{HTML}{E5B400}
\definecolor{milcOcre}{HTML}{B86600}
\definecolor{milcNegro}{HTML}{241816}
\definecolor{milcGris}{HTML}{F7F7F4}
\definecolor{milcLinea}{HTML}{E8E2DD}
\definecolor{milcGradePrimary}{HTML}{FF2D87}
\definecolor{milcGradeDark}{HTML}{9F1B56}
\definecolor{milcGradeSoft}{HTML}{FFE0EE}
\definecolor{milcGradeAccent}{HTML}{CFE7FF}
\definecolor{milcGradeText}{HTML}{FFFFFF}
\color{milcNegro}

\pagestyle{fancy}
\fancyhf{}
% Cabecera de dos lineas: identidad de la guia arriba y, a la derecha y en
% gris, la estacion en curso (\markright desde \milcestacion). Antes de la
% primera estacion la segunda linea va vacia; el \strut mantiene la altura
% para que la linea de arriba no baile entre paginas.
\lhead{\small Tecnología e Informática · Grado 10\\[-1pt]{\footnotesize\strut}}
\rhead{\small Guía 2 · MILC v3\\[-1pt]{\footnotesize\strut\color{milcNegro!65}\rightmark}}
\cfoot{\small\thepage}
\renewcommand{\headrulewidth}{0pt}

% Sobre mostaza (#E5B400) y verde (#58A923) el texto va en milcNegro; sobre el
% resto de la paleta, en blanco. Se usa dentro de fonttitle/fontupper, donde un
% \color posterior manda sobre coltitle.
\newcommand{\milctintasobre}[1]{%
  \ifboolexpr{test{\ifstrequal{#1}{milcMostaza}} or test{\ifstrequal{#1}{milcVerde}}}%
    {\color{milcNegro}}{\color{white}}}

\newtcolorbox{titlebox}[1]{
  enhanced,colback=#1,colframe=#1,arc=7mm,boxrule=0pt,
  left=7mm,right=7mm,top=3mm,bottom=3mm,
  before skip=8pt,after skip=8pt,
  fontupper=\sffamily\bfseries\Large\color{white}
}

\newtcolorbox{softbox}[3]{
  enhanced,breakable,pad at break=3mm,
  colback=#2,colframe=#1,borderline west={4pt}{0pt}{#1},
  arc=2mm,boxrule=.4pt,left=4mm,right=4mm,top=3mm,bottom=3mm,
  before skip=6pt,after skip=8pt,title={#3},coltitle=white,
  colbacktitle=#1,fonttitle=\sffamily\bfseries\milctintasobre{#1},fontupper=\RaggedRight,
  attach boxed title to top left={xshift=3mm,yshift=-2mm},
  boxed title style={arc=2mm, boxrule=0pt}
}

\newtcolorbox{drawbox}[2]{
  enhanced,breakable,pad at break=4mm,
  colback=white,colframe=#1,arc=4mm,boxrule=1pt,
  left=5mm,right=5mm,top=4mm,bottom=4mm,before skip=8pt,after skip=8pt,
  title={#2},coltitle=white,colbacktitle=#1,fonttitle=\sffamily\bfseries\milctintasobre{#1},
  fontupper=\RaggedRight,
  attach boxed title to top left={xshift=4mm,yshift=-2mm},
  boxed title style={arc=3mm, boxrule=0pt}
}

\newtcolorbox{infoband}[1]{
  enhanced,breakable,pad at break=2mm,colback=#1!8,colframe=#1!50,arc=2mm,boxrule=.5pt,
  left=4mm,right=4mm,top=2mm,bottom=2mm,before skip=4pt,after skip=4pt,
  fontupper=\small\RaggedRight
}

% Puente narrativo entre secciones (contrato editorial Conéctate).
% Da continuidad: "Acabas de X. Ahora vas a Y porque Z."
% Visualmente: banda izquierda en color del grado, texto en itálica pequeña.
\newtcolorbox{puentebox}{
  enhanced,breakable,pad at break=2mm,colback=milcGradeSoft,colframe=milcGradeDark,
  borderline west={3pt}{0pt}{milcGradeDark},
  arc=0mm,boxrule=0pt,left=5mm,right=4mm,top=2.5mm,bottom=2.5mm,
  before skip=4pt,after skip=8pt,
  fontupper=\itshape\small\color{milcNegro}\RaggedRight
}
% La flecha va en modo matematico a proposito: Helvetica Neue NO tiene el glifo
% U+2192, asi que \textrightarrow se compilaba a NADA («Missing character: There
% is no → in font Helvetica Neue») y el puente salia sin flecha. En modo
% matematico la flecha viene de la fuente matematica y si se dibuja.
\newcommand{\puente}[1]{%
  \begin{puentebox}%
    \textcolor{milcGradeDark}{$\rightarrow$}\hspace{2mm}#1%
  \end{puentebox}%
}

\newcommand{\linead}[1]{\noindent\color{milcLinea}\rule{#1}{0.5pt}\color{milcNegro}}
\newcommand{\respuesta}[1]{\par\vspace{2mm}\foreach \n in {1,...,#1}{\noindent\color{milcLinea}\rule{\linewidth}{0.5pt}\par\vspace{5mm}}\color{milcNegro}}
\newcommand{\checkbox}{\textcolor{milcGradeDark}{\fbox{\phantom{x}}}\hspace{5pt}}

% ─── Listas aireadas para las guías socioemocionales ────────────────────
% Componer las verificaciones como «\checkbox texto\\» deja el texto que salta
% de línea debajo del cuadrito y apelmaza el bloque. Estos entornos alinean la
% continuación y airean los ítems. Los usa Territorio Interior; cualquier
% familia puede hacerlo.
\newlist{checklista}{itemize}{1}
\setlist[checklista]{
  label=\textcolor{milcGradeDark}{\fbox{\phantom{x}}},
  leftmargin=9mm, labelsep=3.2mm, itemsep=2.4mm,
  topsep=2.5mm, parsep=0pt, partopsep=0pt, before=\RaggedRight
}
\newlist{pasoslista}{enumerate}{1}
\setlist[pasoslista]{
  label=\textcolor{milcGradeDark}{\sffamily\bfseries\arabic*.},
  leftmargin=8mm, labelsep=3mm, itemsep=2.8mm,
  topsep=1mm, parsep=0pt, partopsep=0pt, before=\RaggedRight
}

\newcommand{\phasecard}[4]{
\begin{tcolorbox}[enhanced,colback=#3,colframe=#2,arc=3mm,boxrule=.5pt,height=3.05cm,valign=center,halign=center,left=1.5mm,right=1.5mm,top=2mm,bottom=2mm]
{\sffamily\bfseries\fontsize{22}{24}\selectfont\textcolor{#2}{#1}}\par
{\sffamily\bfseries\small #4}
\end{tcolorbox}
}

% ── Piezas del rediseno 2026 (legibilidad 6.º-11.º) ───────────────────
% Banda de estacion: reemplaza el titlebox macizo. Titulo a la izquierda,
% dato practico (tiempo/modalidad) a la derecha, en una sola linea.
% Nunca huerfana: pide 9 lineas libres (o salta de pagina rellenando con
% \vfill) y prohibe el salto justo despues de la banda. Nueve y no seis:
% la caja partible que sigue necesita ~80pt (titulo adosado, relleno y dos
% lineas) para arrancar en la misma pagina; con menos, tcolorbox la manda
% entera a la siguiente y la banda se queda sola. El penalty 10000 va
% en `after`, ANTES del espacio posterior: un `after skip` (aunque sea 0pt)
% es un \vskip tras la caja, y ese glue es un punto de corte legal que un
% \nopagebreak escrito despues de \end{tcolorbox} ya no cierra. Cada banda
% deja marcador PDF y actualiza la cabecera derecha.
\newcounter{milcestacion}
\newcommand{\milcestacion}[2]{%
\Needspace*{9\baselineskip}%
\stepcounter{milcestacion}%
\pdfbookmark[1]{#2}{milcestacion\themilcestacion}%
\markright{#2}%
\begin{tcolorbox}[enhanced,colback=#1,colframe=#1,arc=2mm,boxrule=0pt,
  left=5mm,right=5mm,top=2.6mm,bottom=2.6mm,before skip=11pt,
  after={\par\nopagebreak\vskip7pt}]
{\sffamily\bfseries\fontsize{15}{18}\selectfont\milctintasobre{#1}#2}
\end{tcolorbox}}

% Tarjeta de paso numerada: sustituye las tablas «Pilar / Aplicacion», que
% eran formato de consulta docente, no de estudiante. El numero va en un
% circulo sobre el borde para que la secuencia se lea de un vistazo.
\newcommand{\milcpaso}[3]{%
\Needspace{4\baselineskip}%
\begin{tcolorbox}[enhanced,colback=white,colframe=milcLinea,arc=1.5mm,boxrule=.9pt,
  left=14mm,right=4mm,top=3mm,bottom=3mm,before skip=4pt,after skip=4pt,
  fontupper=\RaggedRight,
  overlay={\node[anchor=north west,circle,fill=#2,text=white,inner sep=0pt,
    minimum size=7.4mm,font=\sffamily\bfseries\small]
    at ([xshift=3.6mm,yshift=-3.2mm]frame.north west) {#1};}]
#3
\end{tcolorbox}}

% Casilla marcable de tamano util con lapiz (la anterior era \fbox{\phantom{x}},
% de alto variable segun el texto que la seguia).
\newcommand{\milcmarca}{\framebox[3.8mm]{\rule{0pt}{3.4mm}}}

% Fila marcable de lista suelta (checks de la estacion 1).
\newcommand{\milccheckrow}[1]{%
\par\vspace{1.8mm}\noindent
\makebox[6.4mm][l]{\raisebox{-.55mm}{\textcolor{milcNegro}{\milcmarca}}}%
\parbox[t]{\dimexpr\linewidth-6.4mm\relax}{\RaggedRight #1}\par}

% Celda marcable dentro de la rubrica (tabla de autoevaluacion). Misma
% sangria francesa, pero pensada para vivir dentro de una columna X.
\newcommand{\milcopcion}[1]{%
\makebox[5.2mm][l]{\raisebox{-.55mm}{\textcolor{milcNegro}{\milcmarca}}}%
\parbox[t]{\dimexpr\linewidth-5.2mm\relax}{\RaggedRight #1}}

% ── Recursos visuales de la guía ──────────────────────────────────────
% Declarados en el bloque `recursos:` del YAML y emitidos por el builder.
% \guiaFigura   → imagen rasterizada o PDF (foto, carta celeste, montaje).
% \guiaDiagrama → fuente TikZ/circuitikz (.tex) incrustada como vector:
%                 es la vía para circuitos y esquemáticos editables.
% [#1] = texto alternativo (alt) · #2 = ruta relativa al .tex · #3 = pie
% El alt llega al PDF etiquetado (/Alt) y lo lee el lector de pantalla.
\newcommand{\guiaFigura}[3][]{%
\begin{center}
\includegraphics[width=0.88\linewidth,height=0.38\textheight,keepaspectratio,alt={#1}]{#2}\\[1.5mm]
{\footnotesize\itshape\textcolor{milcNegro!75}{#3}}
\end{center}
}

\newcommand{\guiaDiagrama}[2]{%
\begin{center}
\input{#1}\\[1.5mm]
{\footnotesize\itshape\textcolor{milcNegro!75}{#2}}
\end{center}
}

\begin{document}
\thispagestyle{empty}

% ── Portada ───────────────────────────────────────────────────────────
% Antes: un rectangulo de color a pagina completa. Se veia bien en pantalla,
% pero en la fotocopiadora del colegio se comia el toner y salia gris sucio.
% Ahora la banda ocupa el tercio superior y el resto va en blanco.
\begin{tikzpicture}[remember picture,overlay]
  \path[fill=milcGradePrimary]
    (current page.north west) rectangle ([yshift=-10.6cm]current page.north east);

  \node[anchor=north west,text=milcGradeText] at ([xshift=2.0cm,yshift=-1.55cm]current page.north west)
    {\sffamily\bfseries\fontsize{12}{14}\selectfont GRADO};
  \node[anchor=north west,text=milcGradeText,opacity=.85] at ([xshift=2.0cm,yshift=-2.35cm]current page.north west)
    {\sffamily\bfseries\fontsize{8.5}{10}\selectfont SERIE GUÍAS · TECNOLOGÍA E INFORMÁTICA};
  \node[anchor=north east,text=milcGradeText,opacity=.85,align=right] at ([xshift=-2.0cm,yshift=-1.55cm]current page.north east)
    {\sffamily\bfseries\fontsize{9}{12}\selectfont I.E. Sor María Juliana\\Cartago, Valle del Cauca};

  \node[anchor=west,text=milcGradeText] (gradonum) at ([xshift=1.85cm,yshift=-7.1cm]current page.north west)
    {\sffamily\bfseries\fontsize{112}{112}\selectfont 10};
  % El circulito de grado (°) solo aplica a las guias de grado («11°»); en
  % Bebras y Semillero el numero grande es un momento/modulo, no un grado.
  % Se posiciona relativo al nodo `gradonum`, no en coordenadas absolutas.
  \draw[milcGradeText,line width=1.6pt]
    ([xshift=2.0mm,yshift=-4.5mm]gradonum.north east) circle (.15cm);

  \node[anchor=west,text=milcGradeText,text width=11.4cm,align=left] at ([xshift=1.5cm,yshift=0.5cm]gradonum.east)
    {
     {\sffamily\bfseries\fontsize{9}{11}\selectfont\MakeUppercase{Período 1 · Escritura abierta con IA}}\\[3mm]
     {\sffamily\bfseries\fontsize{21}{25}\selectfont\setlength{\baselineskip}{27pt}Concebir tu libro ---\\[4pt]género, tema, audiencia\\[4pt]y escaleta}\\[3.5mm]
     {\sffamily\bfseries\fontsize{15}{18}\selectfont Guía 2-10-TIC}
    };
\end{tikzpicture}

\vspace*{9.4cm}
\vfill

{\sffamily\bfseries\fontsize{8}{10}\selectfont\textcolor{milcGradeDark}{LO QUE TE LLEVAS HOY}}

\begin{tcolorbox}[enhanced,colback=white,colframe=milcNegro,arc=2mm,boxrule=1.6pt,
  left=5mm,right=5mm,top=4mm,bottom=4mm,before skip=3pt,after skip=12pt,
  fontupper=\RaggedRight\fontsize{11}{15.5}\selectfont]
Ficha editorial con género específico, tema, audiencia precisa y propósito, más una escaleta de 8 a 12 capítulos y una justificación de 5 líneas
\end{tcolorbox}

\begin{tcolorbox}[enhanced,colback=milcGris,colframe=milcGris,arc=2mm,boxrule=0pt,
  left=5mm,right=5mm,top=3.5mm,bottom=3.5mm,after skip=12pt]
{\sffamily\bfseries\fontsize{8}{10}\selectfont\textcolor{milcNegro!55}{ESTA GUÍA ES DE}}\\[3mm]
\begin{tabularx}{\linewidth}{X p{3.2cm} p{3.2cm}}
\linead{\linewidth} & \linead{3.0cm} & \linead{3.0cm}\\[-1mm]
{\fontsize{8.5}{10}\selectfont\textcolor{milcNegro!55}{Nombre completo}} &
{\fontsize{8.5}{10}\selectfont\textcolor{milcNegro!55}{Grupo}} &
{\fontsize{8.5}{10}\selectfont\textcolor{milcNegro!55}{Fecha}}\\
\end{tabularx}
\end{tcolorbox}

\vfill

\noindent\textcolor{milcLinea}{\rule{\linewidth}{1pt}}\\[2mm]
{\sffamily\bfseries\fontsize{9.5}{12}\selectfont Autor: Dr. Álvaro Cárdenas Orozco}\\[1mm]
{\sffamily\fontsize{8.5}{11}\selectfont\textcolor{milcNegro!55}{Metodología MILC · apertura ancestral · pensamiento computacional · triángulo Dussel–estoicismo–Floridi}}

\newpage

\section*{Apertura: Concebir tu libro --- género, tema, audiencia y escaleta}

\begin{softbox}{milcGradeDark}{milcGradeSoft}{Saber ancestral}
Los decimeros del Pacífico, en Guapi y en el río Saija, componen décimas glosadas. La forma está decidida antes de la primera palabra. Se declara primero una copla de cuatro versos. Después vienen cuatro décimas, y cada una termina en el verso que le corresponde de esa copla. Son cuarenta y cuatro versos en total. Lo que se va a decir queda anunciado al comienzo, y cada parte cierra regresando a él. La autoría tampoco es fija: cada quien añade o cambia. Y no es folclor. Ulrich Oslender las estudia como discursos ocultos de resistencia de comunidades negras que disputan el reconocimiento de su territorio (Oslender, 2003). La \textbf{cara de exclusión} está en quién sostiene esa memoria, porque la tradición depende de mayores cada vez menos escuchados.\par\smallskip{\footnotesize\textcolor{milcNegro!70}{\textbf{Fuente.} Oslender, U. (2003). «Discursos ocultos de resistencia»: tradición oral y cultura política en comunidades negras de la costa pacífica colombiana. \emph{Revista Colombiana de Antropología, 39}, 203--236.}}
\end{softbox}

\begin{softbox}{milcTurquesa}{milcGris}{Saber contemporáneo}
La \textbf{concepción editorial} es la etapa donde se toman las decisiones que sostienen toda la producción del libro. Antes de pedirle un solo prompt a la IA, se cierran cuatro piezas. El \textbf{género específico}: no basta «ficción», hace falta «novela juvenil de aventura» o «manual práctico paso a paso». El \textbf{tema}: en una o dos frases, qué pregunta explora el libro. La \textbf{audiencia}: edad concreta, contexto y conocimiento previo que se supone. Y el \textbf{propósito}: qué se lleva quien lo lea. A eso se suma la \textbf{escaleta}, la lista preliminar de capítulos con lo que va en cada uno. Sin escaleta, la IA genera capítulos que no conversan entre sí.
\end{softbox}

\begin{softbox}{milcMostaza}{milcGris}{Pregunta-puente}
\textit{¿Por qué un decimero declara los cuatro versos antes de empezar, si podría improvisar sin anunciarlos? ¿Y qué pierde quien le pide un libro a una IA sin haber decidido género, audiencia ni propósito?}
\end{softbox}

\begin{softbox}{milcVerde}{milcGris}{Saber y saber hacer}
Al terminar podrás: (1) \textbf{identificar} entre varios temas propios cuál sostiene un libro de ocho a doce capítulos; (2) \textbf{analizar} tu audiencia concreta con edad, contexto y conocimiento previo, y el género que le corresponde; (3) \textbf{crear} la ficha editorial completa con escaleta y una justificación honesta de por qué este libro merece existir.
\end{softbox}



\small
\begin{tabularx}{\textwidth}{>{\bfseries}p{3.0cm}Y}
\rowcolor{milcGradePrimary!12}Autor & Dr. Álvaro Cárdenas Orozco\\
DBA & Producción editorial mediada por IA con criterio ético y técnico (MEN, T\&I)\\
Referentes & Ley 23 de 1982 · Floridi (ética de la IA generativa) · Dussel (autoría con dignidad)\\
Producto final & Ficha editorial con género específico, tema, audiencia precisa y propósito, más una escaleta de 8 a 12 capítulos y una justificación de 5 líneas\\
\end{tabularx}
\normalsize

\puente{En la sesión 1 aceptaste el rol de editor. Hoy haces la concepción de tu libro: las decisiones que van a sostener las ocho sesiones siguientes. Todo lo que venga después se apoya en lo que decidas hoy.}

\milcestacion{milcGradeDark}{Ruta MILC: así vamos a aprender}
\begin{tabularx}{\textwidth}{>{\centering\arraybackslash}X>{\centering\arraybackslash}X>{\centering\arraybackslash}X>{\centering\arraybackslash}X}
\phasecard{1}{milcVerde}{milcGris}{Escucha\\[-1mm]\small Observo, escucho y pregunto.} &
\phasecard{2}{milcTurquesa}{milcGris}{Entender\\[-1mm]\small Sistematizo y ordeno.} &
\phasecard{3}{milcMagenta}{milcGris}{Hacer\\[-1mm]\small Construyo y aplico.} &
\phasecard{4}{milcVino}{milcGris}{Cierre\\[-1mm]\small Evalúo y reflexiono.}
\end{tabularx}


\puente{Antes de teorizar, vas a hacer una lluvia personal. Cinco temas que te importen, te molesten o te den curiosidad, sin filtrarlos por viabilidad. Solo cosas que te importen de verdad.}

\milcestacion{milcVerde}{Estación 1 de 3 · Escucha}

\begin{softbox}{milcVerde}{milcGris}{Escena para escuchar}
\textbf{Actividad 1 · IDENTIFICA --- Cinco temas propios} (15 min · individual). (1) Anota cinco temas que de verdad te importen, te molesten o te den curiosidad. (2) Que sean tuyos, no los que crees que deberían importarte. (3) Pueden ser amplios o específicos, da igual por ahora. (4) No los censures por difíciles. (5) Escribe al lado de cada uno una línea sobre por qué te importa a ti.
\end{softbox}

{\sffamily\bfseries\fontsize{8}{10}\selectfont\textcolor{milcNegro!55}{MARCA CUANDO LO HAYAS HECHO}}
\milccheckrow{Anoté cinco temas reales míos, no los que quedan bien.}
\milccheckrow{No filtré ninguno por difícil o por poco viable.}
\milccheckrow{Escribí en una línea por qué me importa cada uno.}

\begin{infoband}{milcVerde}


\textbf{Lo que va al cuaderno (Actividad 1 · IDENTIFICA).} \textbf{Título:} «Actividad 1 --- Cinco temas propios». \textbf{Formato:} lista numerada de cinco, con una línea por tema sobre por qué te importa. \textbf{Extensión:} un tercio de página. \textbf{Sabes que terminaste cuando} sientes que cualquiera de los cinco podría ser tu libro.
\end{infoband}

\puente{Ya tienes cinco temas. Ahora vas a entender la \textbf{anatomía de la concepción}: cómo se elige el género, cómo se afina la audiencia, cómo se construye una escaleta y qué errores comete el editor novato.}

\milcestacion{milcTurquesa}{Estación 2 de 3 · Entender}

\begin{softbox}{milcTurquesa}{milcGris}{Lo que vas a entender}
\textbf{Actividad 2 · ANALIZA --- Anatomía de la concepción} (20 min · parejas). (1) Con tu pareja, escriban las cuatro decisiones de concepción con un ejemplo de cada una. (2) Escriban la estructura mínima de una entrada de escaleta. (3) Escriban los cinco errores del editor novato. (4) Cada uno pasa sus cinco temas por los tres filtros y le explica al otro cuál queda de finalista.
\end{softbox}

\milcpaso{1}{milcTurquesa}{Las cuatro decisiones se descomponen así. El \textbf{género} tiene que ser específico: novela juvenil de fantasía, manual práctico de hábitos, cuentos cortos urbanos, divulgación narrativa. El \textbf{tema} cabe en una o dos frases declarativas. La \textbf{audiencia} lleva edad, contexto y nivel previo. El \textbf{propósito} dice qué se lleva quien lee. Cuanto más específico el género, más fácil todo lo demás.}
\milcpaso{2}{milcTurquesa}{La \textbf{escaleta} es la lista preliminar de capítulos. Para un libro de ochenta páginas suelen ser entre ocho y doce, de seis a diez páginas cada uno. Cada entrada lleva número, título tentativo y un resumen de dos o tres líneas de lo que pasa o se explica ahí. La escaleta es el mapa que impide que la IA genere capítulos que no conversan entre sí.}
\milcpaso{3}{milcTurquesa}{Todo se juega en una pregunta: ¿este libro merece existir ahora? La justificación tiene que responder tres cosas. Por qué te importa a ti escribirlo. Por qué le importaría a tu audiencia leerlo. Y qué tiene este libro que no tengan los que ya existen sobre el tema. Si no encuentras las tres, conviene cambiar de tema hoy y no en la sesión 6.}
\milcpaso{4}{milcTurquesa}{Algoritmo: (1) pasa los cinco temas por tres filtros: me importa, cabe en ocho a doce capítulos, le interesaría a alguien que conozco; (2) elige un finalista; (3) declara el género específico; (4) escribe el tema en dos frases; (5) define la audiencia; (6) declara el propósito; (7) construye la escaleta; (8) escribe la justificación.}

\begin{drawbox}{milcTurquesa}{Anatomía de la concepción editorial}
\textbf{Género:} específico, no «ficción» a secas. \\[1mm] \textbf{Tema:} una o dos frases claras. \\[1mm] \textbf{Audiencia:} edad, contexto y nivel previo. \\[1mm] \textbf{Propósito:} qué se lleva quien lee. \\[1mm] \textbf{Escaleta:} de 8 a 12 capítulos con título y resumen. \\[1mm] \textbf{Justificación:} por qué tú, por qué ellos, por qué distinto.
\end{drawbox}

\begin{softbox}{milcMostaza}{milcGris}{Errores comunes y su corrección}
(1) Tema demasiado amplio: «la vida» no cabe en ochenta páginas. (2) Audiencia genérica: «para todos» significa para nadie. (3) Propósito difuso: «compartir mis ideas» no es un propósito. (4) Escaleta inflada: veinte capítulos en ochenta páginas dejan cuatro páginas cada uno y todo queda fragmentado. (5) Un libro que nadie pidió: repetir un manual que ya existe, sin nada propio.
\end{softbox}

\begin{infoband}{milcTurquesa}


\textbf{Lo que va al cuaderno (Actividad 2 · ANALIZA).} \textbf{Título:} «Actividad 2 --- Anatomía de la concepción». \textbf{Formato:} las cuatro decisiones con ejemplo, la estructura de una entrada de escaleta, los cinco errores con el que más temes marcado y el tema finalista. \textbf{Extensión:} media página. \textbf{Sabes que terminaste cuando} podrías auditar la concepción de cualquier libro real.
\end{infoband}

\puente{Tienes el método. Toca aplicarlo. Eliges un tema entre los cinco, cierras las cuatro decisiones, construyes la escaleta y justificas por qué este libro merece escribirse.}

\milcestacion{milcMagenta}{Estación 3 de 3 · Hacer}

\begin{softbox}{milcMagenta}{milcGris}{Cómo vas a construir tu producto}
\textbf{Actividad 3 · CREA --- Mi ficha editorial} (35 min · individual). (1) Declara el género específico en tres a cinco palabras. (2) Escribe el tema en una o dos frases declarativas. (3) Define la audiencia con edad, contexto y conocimiento previo. (4) Declara el propósito en una frase con un verbo de cambio. (5) Construye la escaleta de ocho a doce capítulos, cada uno con título y resumen de dos o tres líneas. (6) Escribe la justificación de cinco líneas y lee toda la ficha en voz alta.
\end{softbox}

\milcpaso{1}{milcMagenta}{Tu ficha se construye en seis piezas: género específico, tema en dos frases, audiencia con edad y contexto, propósito con verbo de cambio, escaleta de ocho a doce capítulos y justificación de cinco líneas. Si falta una, la sesión 3 arranca a ciegas.}
\milcpaso{2}{milcMagenta}{Sigue el patrón «escaleta como mapa, no como cárcel». Es un plan de partida, no un compromiso definitivo. En las sesiones siguientes vas a poder agregar, quitar, fusionar y reordenar capítulos. Pero hoy hace falta un mapa para empezar a generar borradores.}
\milcpaso{3}{milcMagenta}{La justificación es la pieza más honesta de la ficha. Aquí te enfrentas a por qué tú y por qué ahora. «Quiero contar lo que aprendí cuidando a mi abuela porque hay muchos jóvenes en esa situación y nadie les habla» funciona. «Me gusta el tema» no funciona, porque no dice nada de quien va a leer.}
\milcpaso{4}{milcMagenta}{Algoritmo: (1) elige el tema finalista; (2) declara el género; (3) afina tema y audiencia; (4) declara el propósito; (5) escribe la escaleta completa; (6) escribe la justificación; (7) lee la ficha en voz alta y comprueba que se entiende sin que tú la expliques.}

\begin{drawbox}{milcMagenta}{Checklist antes de entregar la ficha}
\checkbox Género específico, no «ficción» a secas. \\[1mm] \checkbox Tema en una o dos frases declarativas. \\[1mm] \checkbox Audiencia con edad y contexto, no «para todos». \\[1mm] \checkbox Propósito con verbo de cambio. \\[1mm] \checkbox Escaleta de 8 a 12 capítulos, cada uno con resumen. \\[1mm] \checkbox Justificación con los tres ingredientes.
\end{drawbox}

\begin{softbox}{milcMagenta}{milcGris}{Plantilla de trabajo}
\textbf{Género.} \textit{Específico, en tres a cinco palabras.} \\[1mm] \textbf{Tema.} \textit{Una o dos frases declarativas.} \\[1mm] \textbf{Audiencia.} \textit{Edad, contexto y nivel previo.} \\[1mm] \textbf{Propósito.} \textit{Un verbo de cambio en una frase.} \\[1mm] \textbf{Escaleta.} \textit{De 8 a 12 capítulos con título y resumen.} \\[1mm] \textbf{Justificación.} \textit{Por qué tú, por qué ellos, por qué distinto.}
\end{softbox}

\begin{infoband}{milcMagenta}


\textbf{Lo que va al cuaderno (Actividad 3 · CREA).} \textbf{Título:} «Actividad 3 --- Ficha editorial de mi libro». \textbf{Formato:} la ficha con sus seis campos y la escaleta completa. \textbf{Extensión:} una página. \textbf{Sabes que terminaste cuando} podrías mostrarle la ficha a alguien de tu casa y entendería qué libro vas a hacer.
\end{infoband}

\puente{Lo que sigue resume \emph{qué} entregas hoy: la ficha editorial, la escaleta y la justificación. El criterio principal es que un compañero que lea tu ficha pueda imaginar el libro completo y decir si lo leería.}

\milcestacion{milcMostaza}{Producto de la sesión}
\textbf{Ficha editorial con los cuatro campos, escaleta de 8 a 12 capítulos y justificación}.
\begin{softbox}{milcMostaza}{milcGris}{Criterios de calidad}
(1) Género específico y no genérico. (2) Tema en una o dos frases declarativas. (3) Audiencia precisa, con edad y contexto. (4) Propósito con verbo de cambio. (5) Escaleta de 8 a 12 capítulos, cada uno con su resumen. (6) Justificación que responde por qué tú, por qué ellos y por qué distinto.
\end{softbox}

\puente{Antes de cerrar, mira la concepción desde las cinco dimensiones humanas. Decidir antes de escribir es respeto por quien va a leer. Improvisar desperdicia el periodo entero.}

\milcestacion{milcVino}{Evaluación liberadora · cinco dimensiones}

\begin{softbox}{milcVino}{milcGris}{Sentido de la evaluación}
Evaluar no es solo revisar una nota. Es reconocer qué aprendí, qué cambió en mí, qué emociones manejé, cómo me relaciono con la ciudad y la comunidad, y qué le dejaré a quien viene detrás.
\end{softbox}

\textbf{Mi reflexión liberadora en cinco dimensiones.} Completa cada frase iniciada:

\small
\begin{tabularx}{\textwidth}{>{\bfseries\RaggedRight\arraybackslash}p{3.4cm}Y>{\RaggedRight\arraybackslash}p{4.6cm}}
\rowcolor{milcVino}\color{white}Dimensión & \color{white}\bfseries Mi respuesta (frame starter) & \color{white}\bfseries Algo que descubrí\\
1. Personal & Lo que más cambió en mí fue \linead{6.5cm} & \linead{4.2cm}\\
2. Emocional & La emoción más fuerte fue \linead{2.5cm} y la regulé \linead{3cm} & \linead{4.2cm}\\
3. Ciudadana & En democracia esto importa porque \linead{6cm} & \linead{4.2cm}\\
4. Local (Cartago) & En mi barrio sería útil para \linead{6.5cm} & \linead{4.2cm}\\
5. Intergeneracional & A mi abuela/abuelo le diría \linead{6.5cm} & \linead{4.2cm}\\
\end{tabularx}
\normalsize

\begin{infoband}{milcVino}
\textbf{Para la próxima sustentación.} Vuelve a esta tabla en seis meses. Verás que las respuestas cambian — no porque lo aprendido cambie, sino porque tú cambias. La evaluación liberadora es una conversación contigo a través del tiempo.
\end{infoband}


\puente{Y para terminar, tres voces sobre decidir antes de hablar. Dussel sobre el registro y el destinatario, Séneca sobre lo mucho que no nutre, Floridi sobre en qué proyecto estamos trabajando.}

\milcestacion{milcGradeDark}{Triángulo de pensamiento · cierre filosófico}

\begin{softbox}{milcVino}{milcGris}{Dussel · lente del nosotros}
\textit{«Hay lengua cotidiana (la de todos los días), lengua de culturas ilustradas, de cultura de masas, lengua de cultura popular… lengua política (que se comprende no por lo que dice sino por lo que calla, contra quién lo dice, cuándo y por qué…).»}

{\footnotesize\itshape\color{milcNegro!70}--- Enrique Dussel · Filosofía de la liberación (1977), §4.2.4.4}

No hay una sola manera de escribir bien: hay registros, y cada uno le habla a alguien. Por eso «para todos» significa para nadie. Cuando defines la edad, el contexto y lo que ya sabe quien va a leer, estás eligiendo un registro. Y esa elección va a decidir cada frase del libro.

\textbf{Si le leyera un capítulo a la persona que dije que era mi audiencia, ¿sonaría escrito para ella?}
\end{softbox}
\linead{\linewidth}\\[1mm]
\linead{\linewidth}

\begin{softbox}{milcMostaza}{milcGris}{Estoicismo · Séneca · Cartas a Lucilio, 2 (c. 64 d.C.) · cuidado interior}
\textit{«Empalagarse con muchas cosas es lo propio de los estómagos hastiados. Lo mucho y lo muy diverso, no nutre: contamina.»}

{\footnotesize\itshape\color{milcNegro!70}--- Séneca · Cartas a Lucilio, 2 (c. 64 d.C.)}

Tienes cinco temas y vas a escribir uno. Meter los cinco en el mismo libro es la tentación, y es la manera más segura de que ninguno quede dicho. El decimero declara cuatro versos, no cuarenta. La forma existe para obligar a elegir.

\textbf{De los cinco temas, ¿elegí uno de verdad o metí pedazos de los otros cuatro en la escaleta?}
\end{softbox}
\linead{\linewidth}\\[1mm]
\linead{\linewidth}

\begin{softbox}{milcTurquesa}{milcGris}{Floridi · lente de la infoesfera}
\textit{«Una de las preguntas políticas apremiantes que enfrentamos en las sociedades de la información avanzadas es: ¿en qué clase de proyecto humano estamos trabajando? (trad. propia)»}

{\footnotesize\itshape\color{milcNegro!70}--- Luciano Floridi · Commentary on the Onlife Manifesto (2015), § 3.1}

La justificación de cinco líneas es tu versión pequeña de esa pregunta. Producir texto ya no cuesta nada, así que la pregunta dejó de ser si puedes hacer un libro. La pregunta es para qué vas a sumar uno más, y a quién le hace falta.

\textbf{Si este libro no existiera, ¿qué se perdería, y quién lo notaría?}
\end{softbox}
\linead{\linewidth}\\[1mm]
\linead{\linewidth}

\begin{infoband}{milcNegro}
\textbf{Nota para el docente.} Las tres son citas textuales y llevan su referencia debajo. Puedes remitir a la obra en clase.
\end{infoband}

\milcestacion{milcVino}{Mi autoevaluación · marca una casilla por fila}

{\small
\setlength{\extrarowheight}{2.2pt}
\begin{tabularx}{\textwidth}{>{\RaggedRight\arraybackslash\bfseries}p{3.0cm}YYY}
\rowcolor{milcVino}\color{white}\bfseries Criterio & \color{white}\bfseries Lo logré & \color{white}\bfseries En proceso & \color{white}\bfseries Necesito apoyo\\[.6mm]
Comprensión del tema & \milcopcion{Explico las ideas con mis palabras.} & \milcopcion{Reconozco las ideas, falta precisión.} & \milcopcion{Necesito repasar lo básico.}\\[.8mm]
\rowcolor{milcGris}Pensamiento computacional & \milcopcion{Usé los pilares al armar mi producto.} & \milcopcion{Usé algunos pilares.} & \milcopcion{Necesito releer la estación 2.}\\[.8mm]
Aplicación y producto & \milcopcion{Mi producto cumple los criterios.} & \milcopcion{Está armado, con ajustes pendientes.} & \milcopcion{Aún está en construcción.}\\[.8mm]
\rowcolor{milcGris}Reflexión 5 dimensiones & \milcopcion{Respondí cada una con honestidad.} & \milcopcion{Respondí algunas con detalle.} & \milcopcion{Mis respuestas son muy generales.}\\[.8mm]
Triángulo de pensamiento & \milcopcion{Conecté las 3 voces con mi trabajo.} & \milcopcion{Conecté al menos una.} & \milcopcion{No las relaciono con mi proceso.}\\[.8mm]
\end{tabularx}}

\vspace{3mm}
\textbf{Hoy entendí que:}\\[1mm]
\linead{\linewidth}\\[1mm]
\linead{\linewidth}

\textbf{Mi compromiso para la próxima sesión es:}\\[1mm]
\linead{\linewidth}\\[1mm]
\linead{\linewidth}

\begin{softbox}{milcMostaza}{milcGris}{Cita de despedida · Marco Aurelio}
\textit{``El obstáculo en el camino se vuelve el camino. Lo que estorba al camino se vuelve camino.''}\\[1mm]
Cada error de hoy, bien observado, se convierte en pieza del camino que pisarás mañana.
\end{softbox}

\small
\begin{tabularx}{\textwidth}{>{\bfseries}p{3.6cm}Y}
\rowcolor{milcGradeDark!12}Nombre completo: & \linead{10cm}\\
Fecha: & \linead{4cm} \quad Curso/grupo: \linead{4cm}\\
Mi firma: & \linead{9cm}\\
\end{tabularx}
\normalsize

\begin{infoband}{milcGradeDark}
\textbf{Para la próxima sesión.} Llega con la ficha editorial completa y la escaleta de ocho a doce capítulos. En la sesión 3 vas a aprender a escribir prompts profesionales para generar los primeros borradores. La concepción de hoy es el plano de todo eso.
\end{infoband}

\end{document}
